% GENERATED FILE. Edit canonical lessons, then run npm run generate:books.
% canonical-source-hash: fnv1a64:d5d1aeef6adda08c
% canonical-lessons: LA-C339-proelium, LA-C339-acies, LA-C339-agmen, LA-C339-incursio, LA-C339-dimico

\chapter{Proelium}
\label{ch:la-proelium}
\hlchaptermodality{339}

\begin{chapteropening}
\textbf{By the end of this chapter:} \emph{I can tell the story of a battle: the battle, the battle line, the army on the march, a raid and to fight it out.}

Complete the last lesson of chapter 339: I can tell the story of a battle: the battle, the battle line, the army on the march, a raid and to fight it out.
\end{chapteropening}

\section[proelium, proeliī]{proelium, proeliī — a battle}

\label{lesson:LA-C339-proelium}



\begin{quote}
Before the new one: say the Latin for a standard-bearer, then the Latin for a tribune.
\end{quote}

\subsection*{You'll want to know: proelium, proeliī}
\textbf{proelium, proeliī} — \textquotedblleft{}a battle\textquotedblright{}.

\begin{quote}
Proelium ab hōrā prīmā ad vesperum dūrābat.
\end{quote}

\begin{quote}
\emph{The battle lasted from the first hour until evening.}
\end{quote}

Neuter. A single battle; \emph{bellum} (taught) is the whole war. \emph{Proelium committere}: to join battle.

\begin{cousinweb}
\textbf{proelium} is a word whose origin Lewis \& Short call doubtful; they suggest, with hesitation, a link with \emph{duellum}, the old form of \emph{bellum}.
\end{cousinweb}

\begin{grammarlens}[title={What you've built}]
One more word for telling, explaining and describing at B1.
\end{grammarlens}

\subsection*{Your turn}
\begin{itemize}
\raggedright
  \item \emph{Say it:} \emph{proelium, proeliī}
  \item \emph{Say it:} \emph{proelium, proeliī}, once more
  \item \emph{Say it:} a sentence of your own with \emph{proelium, proeliī}
  \item \emph{Recall:} say the Latin for a standard-bearer, then the Latin for a tribune, then say \emph{proelium, proeliī} again
\end{itemize}

\begin{tcolorbox}[breakable,colback=teal!4,colframe=teal!35!black,title={Before you move on}]
What does proelium, proeliī mean? (\textquotedblleft{}A battle\textquotedblright{}.) Say it once more.
\end{tcolorbox}
\section[aciēs, aciēī]{aciēs, aciēī — a battle line; a sharp edge}

\label{lesson:LA-C339-acies}



\begin{quote}
Before the new one: say the Latin for a tribune, then the Latin for a battle.
\end{quote}

\subsection*{You'll want to know: aciēs, aciēī}
\textbf{aciēs, aciēī} — \textquotedblleft{}a battle line; a sharp edge\textquotedblright{}.

\begin{quote}
Dux aciem ante castra posuit.
\end{quote}

\begin{quote}
\emph{The leader drew up the battle line in front of the camp.}
\end{quote}

Feminine, 5th declension, like \emph{diēs}. First a sharp edge, then the sharp front of an army.

\begin{cousinweb}
\textbf{aciēs} is inherited from a Proto-Indo-European root meaning \textquotedblleft{}sharp\textquotedblright{} (EtymDB): the edge of a blade, then the front edge of an army. English \emph{acid} comes from the same root.
\end{cousinweb}

\begin{grammarlens}[title={What you've built}]
One more word for telling, explaining and describing at B1.
\end{grammarlens}

\subsection*{Your turn}
\begin{itemize}
\raggedright
  \item \emph{Say it:} \emph{aciēs, aciēī}
  \item \emph{Say it:} \emph{aciēs, aciēī}, once more
  \item \emph{Say it:} a sentence of your own with \emph{aciēs, aciēī}
  \item \emph{Recall:} say the Latin for a tribune, then the Latin for a battle, then say \emph{aciēs, aciēī} again
\end{itemize}

\begin{tcolorbox}[breakable,colback=teal!4,colframe=teal!35!black,title={Before you move on}]
What does aciēs, aciēī mean? (\textquotedblleft{}A battle line; a sharp edge\textquotedblright{}.) Say it once more.
\end{tcolorbox}
\section[agmen, agminis]{agmen, agminis — an army on the march; a crowd}

\label{lesson:LA-C339-agmen}



\begin{quote}
Before the new one: say the Latin for a battle, then the Latin for a battle line.
\end{quote}

\subsection*{You'll want to know: agmen, agminis}
\textbf{agmen, agminis} — \textquotedblleft{}an army on the march; a crowd\textquotedblright{}.

\begin{quote}
Agmen longum per vallem prōcēdēbat.
\end{quote}

\begin{quote}
\emph{A long line of soldiers was moving forward through the valley.}
\end{quote}

Neuter, 3rd declension: a body in motion. \emph{Novissimum agmen}: the rearguard.

\begin{cousinweb}
\textbf{agmen} is built on \emph{agō}, \textquotedblleft{}to drive, to lead\textquotedblright{} (EtymDB), with the ending \emph{-men}: a train in motion.
\end{cousinweb}

\begin{grammarlens}[title={What you've built}]
One more word for telling, explaining and describing at B1.
\end{grammarlens}

\subsection*{Your turn}
\begin{itemize}
\raggedright
  \item \emph{Say it:} \emph{agmen, agminis}
  \item \emph{Say it:} \emph{agmen, agminis}, once more
  \item \emph{Say it:} a sentence of your own with \emph{agmen, agminis}
  \item \emph{Recall:} say the Latin for a battle, then the Latin for a battle line, then say \emph{agmen, agminis} again
\end{itemize}

\begin{tcolorbox}[breakable,colback=teal!4,colframe=teal!35!black,title={Before you move on}]
What does agmen, agminis mean? (\textquotedblleft{}An army on the march; a crowd\textquotedblright{}.) Say it once more.
\end{tcolorbox}
\section[incursiō, incursiōnis]{incursiō, incursiōnis — a raid, an attack}

\label{lesson:LA-C339-incursio}



\begin{quote}
Before the new one: say the Latin for a battle line, then the Latin for an army on the march.
\end{quote}

\subsection*{You'll want to know: incursiō, incursiōnis}
\textbf{incursiō, incursiōnis} — \textquotedblleft{}a raid, an attack\textquotedblright{}.

\begin{quote}
Hostēs incursiōnem in agrōs nostrōs fēcērunt.
\end{quote}

\begin{quote}
\emph{The enemy made a raid on our fields.}
\end{quote}

Feminine, 3rd declension, from \emph{incurrō}, 'to run against'. English \emph{incursion}.

\begin{cousinweb}
\textbf{incursiō} is built on \emph{incurrō}, \textquotedblleft{}to run against\textquotedblright{}, from \emph{in} and \emph{currō}, \textquotedblleft{}to run\textquotedblright{} (Lewis \& Short). English \emph{incursion} comes from it.
\end{cousinweb}

\begin{grammarlens}[title={What you've built}]
One more word for telling, explaining and describing at B1.
\end{grammarlens}

\subsection*{Your turn}
\begin{itemize}
\raggedright
  \item \emph{Say it:} \emph{incursiō, incursiōnis}
  \item \emph{Say it:} \emph{incursiō, incursiōnis}, once more
  \item \emph{Say it:} a sentence of your own with \emph{incursiō, incursiōnis}
  \item \emph{Recall:} say the Latin for a battle line, then the Latin for an army on the march, then say \emph{incursiō, incursiōnis} again
\end{itemize}

\begin{tcolorbox}[breakable,colback=teal!4,colframe=teal!35!black,title={Before you move on}]
What does incursiō, incursiōnis mean? (\textquotedblleft{}A raid, an attack\textquotedblright{}.) Say it once more.
\end{tcolorbox}
\section[dīmicō, dīmicāre]{dīmicō, dīmicāre — to fight it out, to struggle}

\label{lesson:LA-C339-dimico}



\begin{quote}
Before the new one: say the Latin for an army on the march, then the Latin for a raid.
\end{quote}

\subsection*{You'll want to know: dīmicō, dīmicāre}
\textbf{dīmicō, dīmicāre} — \textquotedblleft{}to fight it out, to struggle\textquotedblright{}.

\begin{quote}
Ambō exercitūs ācriter dīmicāvērunt.
\end{quote}

\begin{quote}
\emph{Both armies fought it out fiercely.}
\end{quote}

\emph{dīmicō, dīmicāre, dīmicāvī, dīmicātum}. A hard, deciding fight.

\begin{cousinweb}
\textbf{dīmicō} is built on \emph{micō}, \textquotedblleft{}to flash, to quiver\textquotedblright{}, with \emph{dis-} (EtymDB): perhaps first the flashing of swords to and fro.
\end{cousinweb}

\begin{grammarlens}[title={What you've built}]
One more word for telling, explaining and describing at B1.
\end{grammarlens}

\subsection*{Your turn}
\begin{itemize}
\raggedright
  \item \emph{Say it:} \emph{dīmicō, dīmicāre}
  \item \emph{Say it:} \emph{dīmicō, dīmicāre}, once more
  \item \emph{Say it:} a sentence of your own with \emph{dīmicō, dīmicāre}
  \item \emph{Recall:} say the Latin for an army on the march, then the Latin for a raid, then say \emph{dīmicō, dīmicāre} again
\end{itemize}

\begin{tcolorbox}[breakable,colback=teal!4,colframe=teal!35!black,title={Before you move on}]
What does dīmicō, dīmicāre mean? (\textquotedblleft{}To fight it out, to struggle\textquotedblright{}.) Say it once more.
\end{tcolorbox}
